\section{The Vacuum Baseline}
\label{sec:the_test_case}
From living cells and neural networks to sprawling cities and economic ecologies, Complex Adaptive Systems (CAS) research has historically yielded rich qualitative taxonomies. IE introduces a falsifiable paradigm shift: by applying the exact constraints of the six persistence pillars to CAS, open-ended qualitative observation is replaced by strict structural requirements.

While macroscopic domains (biological, ecological, economic) must manage the six pillars probabilistically within chaotic, open environments, we require an absolute baseline to test the architecture's universal validity. We postulate that the fundamental quantum vacuum is the simplest, most constrained Complex Adaptive System (CAS) in existence. By defining the quantum vacuum as a persistent system in a state of Quiescent Equilibrium, operating with zero external flux and strictly subject to the six IE pillars, we establish an absolute boundary condition.

As the closed, persistent substrate from which all observable structures emerge, the vacuum admits no deeper system beneath it. It cannot obtain configuration information, boundary conditions, synchronization signals, or memory from an external source. We formalize this closed-system condition as:
\begin{quote}
\textbf{The Closure Constraint:} The substrate of the vacuum must be entirely self-defining. Its geometry must arise exclusively from its own architectural requirements.
\end{quote}
Any non-zero selection entropy would require an external memory to resolve the degeneracy, violating self-specification.
\label{subsec:closure}

Subjecting the vacuum to this constraint creates the most constrained test case possible: if persistence requirements alone cannot derive its mathematical structure, the framework is not universal. This paper strictly bounds its scope to the static geometric baseline; the continuous dynamic Effective Field Theory (EFT) Lagrangian is left to future work.
